\documentclass[11pt]{article}

% ============================================================
% Packages
% ============================================================
\usepackage[a4paper,margin=1in]{geometry}
\usepackage[T1]{fontenc}
\usepackage{lmodern}
\usepackage{microtype}
\usepackage{amsmath,amssymb,amsthm,mathtools}
\usepackage{bm}
\usepackage{xcolor}
\usepackage[colorlinks=true,linkcolor=blue!55!black,urlcolor=blue!55!black,citecolor=blue!55!black]{hyperref}
\usepackage[nameinlink,capitalize]{cleveref}

% ============================================================
% Theorem environments
% ============================================================
\theoremstyle{plain}
\newtheorem{theorem}{Theorem}[section]
\newtheorem{proposition}[theorem]{Proposition}
\newtheorem{lemma}[theorem]{Lemma}
\newtheorem{corollary}[theorem]{Corollary}

\theoremstyle{definition}
\newtheorem{definition}[theorem]{Definition}
\newtheorem{assumption}[theorem]{Assumption}

\theoremstyle{remark}
\newtheorem{remark}[theorem]{Remark}

% ============================================================
% Macros
% ============================================================
\newcommand{\R}{\mathbb{R}}
\newcommand{\E}{\mathbb{E}}
\newcommand{\Prob}{\mathbb{P}}
\newcommand{\N}{\mathcal{N}}
\newcommand{\Sym}{\mathbb{S}}
\newcommand{\Tr}{\operatorname{Tr}}
\newcommand{\diag}{\operatorname{diag}}
\newcommand{\op}{\mathrm{op}}
\newcommand{\inj}{\mathrm{inj}}
\newcommand{\Lip}{\operatorname{Lip}}
\newcommand{\Id}{I}
\newcommand{\dd}{\mathrm{d}}
\newcommand{\1}{\mathbf{1}}
\newcommand{\norm}[1]{\left\lVert #1\right\rVert}
\newcommand{\abs}[1]{\left\lvert #1\right\rvert}
\newcommand{\inner}[2]{\left\langle #1,#2\right\rangle}

% ============================================================
% Title
% ============================================================
\title{A Convex-Ridge Counterexample to Dimension-Free Logarithmic Local Rates for Gaussian Fisher--Rao Variational Inference}
\author{}
\date{15 July 2026}

\begin{document}
\maketitle

\begin{abstract}
We construct a sequence of strongly convex potentials satisfying exact Gaussian
equilibrium whitening and a dimension-free global Lipschitz bound on the Hessian,
for which the first-Hermite Hessian coupling grows polynomially in the condition
number and the local Gaussian Fisher--Rao spectral gap is at most of order
$\kappa^{-1/2}$. The construction is a high-dimensional convex ridge whose
orientation is weakly tilted toward one distinguished coordinate. Its averaged
Hessian is tuned to be exactly isotropic. A distributional Codazzi identity,
equivalently a direct radial Gaussian integration-by-parts identity, exposes a
collective radial fluctuation that produces a large first-Hermite coupling. The
same coupling produces an explicit nearly-null
Rayleigh direction for the packed local generator. Consequently, neither a
dimension-free Codazzi--Chang logarithmic estimate for the first-Hermite coupling
nor a dimension-free $(1+\log\kappa)^{-1}$ lower bound on the local convergence
rate can hold under a globally Lipschitz Hessian assumption alone.
\end{abstract}

\tableofcontents

% ============================================================
\section{Setting and main result}
% ============================================================

Let $n\geq1$, let $Z\sim\N(0,\Id_n)$, and let
$V:\R^n\to\R$ be a $C^2$ potential. We consider the assumptions
\begin{align}
    \alpha\Id_n
    \preceq
    \nabla^2V(x)
    \preceq
    \beta\Id_n,
    \qquad x\in\R^n,
    \qquad
    0<\alpha\leq1\leq\beta,
    \label{eq:H1}\\
    \E[\nabla V(Z)]=0,
    \qquad
    \E[\nabla^2V(Z)]=\Id_n.
    \label{eq:H2}
\end{align}
We write
\[
    \kappa:=\frac{\beta}{\alpha},
    \qquad
    W(z):=\nabla^2V(z),
    \qquad
    A(z):=W(z)-\alpha\Id_n,
    \qquad
    a:=1-\alpha,
    \qquad
    L:=\beta-\alpha.
\]
Then
\[
    0\preceq A(z)\preceq L\Id_n,
    \qquad
    \E[A(Z)]=a\Id_n.
\]

For $X\in\Sym(n)$, define the first-Hermite contraction
\begin{equation}
    T[X]
    :=
    \E\!\left[\Tr(XW(Z))Z\right]
    \in\R^n.
    \label{eq:T-def}
\end{equation}
For any $w\in\R^n$ and $X\in\Sym(n)$, the scalar $\inner{w}{T[X]}$ equals
\[
    \E[(w^\top Z)\Tr(XW(Z))]
    =\inner{X}{\E[(w^\top Z)W(Z)]}_{\mathrm F},
\]
so the adjoint of $T$ (with respect to the Frobenius and Euclidean inner
products) is
\begin{equation}
    T^*w
    =
    \E\!\left[(w^\top Z)W(Z)\right],
    \qquad w\in\R^n.
    \label{eq:Tstar-def}
\end{equation}
Hence the first-Hermite coupling norm is
\begin{equation}
    \tau_H
    :=
    \sup_{\norm{w}=1}
    \norm{\E[(w^\top Z)W(Z)]}_{\mathrm F}
    =
    \norm{T^*}_{\R^n\to\Sym(n)}
    =
    \norm{T}_{\Sym(n)\to\R^n}.
    \label{eq:tau-def}
\end{equation}

We also define the second-Hermite operator
\begin{equation}
    H[X]
    :=
    \E\!\left[
        \bigl(Z^\top XZ-\Tr X\bigr)\bigl(W(Z)-\Id_n\bigr)
    \right],
    \qquad X\in\Sym(n).
    \label{eq:H-def}
\end{equation}
That $H$ is self-adjoint on $\Sym(n)$ is proved in
\Cref{lem:second-Hermite} below; it is used implicitly in the definition of
$L_\star$ and explicitly in the variational argument of \Cref{sec:gap}.
At the whitened equilibrium, the packed self-adjoint local generator is
\begin{equation}
    L_\star
    :=
    \begin{pmatrix}
        \Id_n & T/\sqrt2\\[2mm]
        T^*/\sqrt2 & \Id+H/2
    \end{pmatrix}
    \quad\text{on}\quad
    \R^n\oplus\Sym(n),
    \label{eq:Lstar-def}
\end{equation}
where the packed norm is
\[
    \norm{(v,X)}_{\mathrm{pack}}^2
    :=
    \norm v^2+\norm X_{\mathrm F}^2.
\]
The off-diagonal blocks of $L_\star$ are mutually adjoint by
\eqref{eq:Tstar-def}, and the diagonal blocks are self-adjoint, so $L_\star$
is self-adjoint and its spectrum is real. The local convergence rate is
\[
    \gamma_{\mathrm{loc}}
    :=
    \lambda_{\min}(L_\star).
\]

The optional regularity assumption considered below is
\begin{equation}
    \norm{\nabla^2V(x)-\nabla^2V(y)}_{\op}
    \leq
    L_H\norm{x-y},
    \qquad x,y\in\R^n.
    \label{eq:H3}
\end{equation}

\begin{theorem}[Failure of both dimension-free logarithmic conclusions]
\label{thm:main}
There are numerical constants $K_0,c_0,C_0,C_1,C_1'>0$ with the following
property. For every integer $m$ with $(8m)^{1/4}\geq K_0$, set
\[
    K:=(8m)^{1/4},
    \qquad
    n:=m+1,
\]
so that $K\to\infty$ as $m\to\infty$. There is a potential $V:\R^n\to\R$
satisfying \eqref{eq:H1}--\eqref{eq:H2} such that
\begin{align}
    c_0K^2
    \leq
    \kappa
    \leq
    C_0K^2,
    \label{eq:kappa-main}\\
    \tau_H^2
    \geq
    \frac{K}{2}-C_1,
    \label{eq:tau-main}\\
    \gamma_{\mathrm{loc}}
    \leq
    \frac{C_1}{K}
    \leq
    \frac{C_1'}{\sqrt\kappa}.
    \label{eq:gap-main}
\end{align}
Moreover, every potential in the sequence satisfies \eqref{eq:H3} with the
same admissible constant
\[
    L_H=64.
\]

Consequently:

\begin{enumerate}
\item There is no absolute constant $C$ such that
\[
    \tau_H^2
    \leq
    C a^2\left(1+\log\frac{eL}{a}\right)
\]
holds in every dimension under \eqref{eq:H1}--\eqref{eq:H2}.

\item There is no positive function $c(L_H)$ such that
\[
    \gamma_{\mathrm{loc}}
    \geq
    \frac{c(L_H)}{1+\log\kappa}
\]
holds in every dimension under \eqref{eq:H1}--\eqref{eq:H3}.
\end{enumerate}
\end{theorem}

The remainder of the note proves \Cref{thm:main}. All constants implicit in
$O(\cdot)$ are numerical and independent of the dimension and of $K$. The
assembly of the individual lemmas into the theorem is recorded at the end of
\Cref{sec:gap}.

% ============================================================
\section{Construction and exact whitening}
% ============================================================

\subsection{A one-dimensional convex profile}

Let $m\in\mathbb N$ be large and define
\begin{equation}
    K:=(8m)^{1/4},
    \qquad
    n:=m+1.
    \label{eq:K-m}
\end{equation}
Thus
\[
    m=\frac{K^4}{8}.
\]
Write $z=(t,y)\in\R\times\R^m$, and set
\begin{equation}
    r:=\norm y,
    \qquad
    \rho(y):=\sqrt{m+r^2},
    \qquad
    \rho_0:=\sqrt{2m}=\frac{K^2}{2},
    \qquad
    c:=\rho_0+\frac12.
    \label{eq:rho-def}
\end{equation}

Define a convex $C^2$ function $F:\R\to\R$ by
\begin{equation}
    F''(x):=(K-\abs{x-c})_+,
    \qquad
    F'(-\infty)=0.
    \label{eq:F-def}
\end{equation}
An additive constant in $F$ is irrelevant and may be fixed arbitrarily. The
profile satisfies
\begin{equation}
    0\leq F''\leq K,
    \qquad
    0\leq F'\leq K^2,
    \qquad
    \Lip(F'')\leq1.
    \label{eq:F-basic}
\end{equation}
Indeed $F''$ is supported in $[c-K,c+K]$ with $\int_\R F''=K^2$, whence the
bound on $F'=\int_{-\infty}^{\,\cdot}F''$; and $\abs{F'''}\leq1$ almost
everywhere, since $F'''(x)=-\operatorname{sgn}(x-c)$ on
$0<\abs{x-c}<K$ and $F'''=0$ elsewhere. For $s\geq0$, define
\begin{equation}
    U_s(t,y):=F(\rho(y)+st),
    \qquad
    x:=\rho(y)+st,
    \qquad
    g:=\nabla x=\left(s,\frac{y}{\rho}\right).
    \label{eq:Us-def}
\end{equation}
Since
\[
    D^2\rho(y)
    =
    \frac{\Id_m}{\rho}
    -
    \frac{yy^\top}{\rho^3},
\]
the Hessian of $U_s$ is
\begin{equation}
    A_s:=D^2U_s
    =
    F''(x)gg^\top
    +
    F'(x)
    \begin{pmatrix}
        0&0\\
        0&D^2\rho(y)
    \end{pmatrix}.
    \label{eq:As}
\end{equation}
Both summands are positive semidefinite: $F''\geq0$, $F'\geq0$, and
$D^2\rho\succeq0$ because $\rho$ is concave-free here, having eigenvalue
$1/\rho>0$ on $y^\perp$ and $m/\rho^3\geq0$ along $y$.

\subsection{Radial concentration}

Let $(T,Y)\sim\N(0,\Id_1)\otimes\N(0,\Id_m)$ and continue to write
$r=\norm Y$ and $\rho=\sqrt{m+r^2}$. Define
\[
    \xi:=\rho-\rho_0.
\]

\begin{lemma}[Dimension-free radial fluctuations]
\label{lem:xi}
For every fixed $p<\infty$, there is a numerical constant $C_p$ such that
\begin{equation}
    \norm{\xi}_{L^p}\leq C_p.
    \label{eq:xi-Lp}
\end{equation}
Moreover,
\begin{equation}
    \E\xi
    =
    -\frac{\E\xi^2}{2\rho_0}
    =
    O(K^{-2}).
    \label{eq:xi-mean}
\end{equation}
Uniformly for $0\leq s\leq2K^{-1/2}$, the variable
\[
    d:=\xi+sT-\frac12
\]
has dimension-free sub-Gaussian tails. In particular, for every fixed
$q<\infty$, there are numerical constants $c_q,C_q>0$ such that
\begin{equation}
    \sup_{0\leq s\leq2K^{-1/2}}
    \E\!\left[
        (1+\abs d)^q\1_{\{\abs d>K\}}
    \right]
    \leq
    C_qe^{-c_qK^2}.
    \label{eq:d-tail}
\end{equation}
\end{lemma}

\begin{proof}
The identity
\[
    \xi
    =
    \frac{r^2-m}{\rho+\rho_0}
\]
and the lower bound $\rho+\rho_0\geq\rho_0\asymp\sqrt m$ imply
\[
    \norm\xi_{L^p}
    \leq
    \frac{\norm{r^2-m}_{L^p}}{\rho_0}.
\]
The centered moments of a $\chi_m^2$ random variable satisfy
$\norm{r^2-m}_{L^p}\leq C_p\sqrt m$ (for instance by writing
$r^2-m=\sum_{i=1}^m(Y_i^2-1)$ and applying Rosenthal's inequality), which
proves \eqref{eq:xi-Lp}. Since
\[
    \rho^2-\rho_0^2
    =
    2\rho_0\xi+\xi^2
\]
and $\E(\rho^2-\rho_0^2)=\E(r^2-m)=0$, taking expectations gives
$2\rho_0\E\xi+\E\xi^2=0$, i.e.\ \eqref{eq:xi-mean}; the final bound uses
$\E\xi^2\leq C$ from \eqref{eq:xi-Lp} and $\rho_0=K^2/2$.

The map $y\mapsto\sqrt{m+\norm y^2}$ is $1$-Lipschitz, so by Gaussian
concentration $\rho-\E\rho$ is sub-Gaussian with an absolute variance proxy;
combined with $\abs{\E\rho-\rho_0}=\abs{\E\xi}=O(K^{-2})$, the centered
variable $\xi$ is uniformly sub-Gaussian. The independent variable $sT$ is
sub-Gaussian with variance proxy $s^2\leq4/K$. Hence $d=\xi+sT-\tfrac12$ is
uniformly sub-Gaussian with an absolute proxy, so
$\Prob(\abs d>K)\leq2e^{-c K^2}$; the weighted tail estimate
\eqref{eq:d-tail} follows by Cauchy--Schwarz against the polynomial weight,
whose $L^2$ norm grows at most polynomially in $K$.
\end{proof}

\subsection{The averaged Hessian and the tuning equation}

\begin{lemma}[Diagonal average]
\label{lem:mean-As}
For every $s\geq0$,
\begin{equation}
    \E A_s
    =
    \diag\bigl(s^2\overline K(s),G(s)\Id_m\bigr),
    \label{eq:mean-As}
\end{equation}
where
\begin{align}
    \overline K(s)
    &:=
    \E F''(\rho+sT),
    \label{eq:barK}\\
    G(s)
    &:=
    \frac1m
    \E\!\left[
        F''(\rho+sT)\frac{r^2}{\rho^2}
        +
        F'(\rho+sT)
        \left(
            \frac m\rho-\frac{r^2}{\rho^3}
        \right)
    \right].
    \label{eq:G-def}
\end{align}
\end{lemma}

\begin{proof}
Write $\E A_s=\E[F''gg^\top]+\E[F'\,\diag(0,D^2\rho)]$ using
\eqref{eq:As}, and recall $g=(s,y/\rho)$. The $(t,t)$ entry is
$\E[F''s^2]=s^2\overline K(s)$. The $(t,y_i)$ entry is
$\E[F''sy_i/\rho]$; since $F''=F''(\rho+sT)$ and $\rho$ depend on $y$ only
through $r=\norm y$, the integrand is odd under the reflection
$y_i\mapsto-y_i$, which preserves the Gaussian law, so the entry vanishes.
For $i\neq j$ the $(y_i,y_j)$ entry is
$\E[F''y_iy_j/\rho^2-F'y_iy_j/\rho^3]=0$ by the same reflection. The diagonal
entries $(y_i,y_i)$ are all equal by exchangeability of the coordinates of
$Y$; their common value equals the average over $i=1,\dots,m$, which by
\eqref{eq:As} and $\sum_i y_i^2=r^2$ is exactly $G(s)$.
\end{proof}

\begin{lemma}[Uniform averaged-profile estimates]
\label{lem:profile-est}
Uniformly for $0\leq s\leq2K^{-1/2}$,
\begin{align}
    \overline K(s)
    &=
    K+O(1),
    \label{eq:barK-est}\\
    \E F'(\rho+sT)
    &=
    \rho_0-\frac K2+O(1),
    \label{eq:Fprime-est}\\
    G(s)
    &=
    1-\frac1K+O(K^{-2}).
    \label{eq:G-est}
\end{align}
\end{lemma}

\begin{proof}
With $d=\xi+sT-\tfrac12$, one has $\rho+sT=c+d$. On the event
$\abs d\leq K$, the definition of $F$ gives the exact formulas
\begin{align}
    F''(c+d)&=K-\abs d,
    \label{eq:Fpp-local}\\
    F'(c+d)&=\frac{K^2}{2}+Kd-\frac12d\abs d.
    \label{eq:Fp-local}
\end{align}
(For \eqref{eq:Fp-local}, integrate \eqref{eq:Fpp-local} from $c$, using
$F'(c)=K^2/2$, which holds because $F''$ is symmetric about $c$ with total
mass $K^2$.) By \Cref{lem:xi}, $\E\abs d=O(1)$ and $\E[d\abs d]=O(1)$
uniformly in $s$. The contribution from $\{\abs d>K\}$ is exponentially small
by \eqref{eq:d-tail}, even after multiplication by any fixed polynomial in
$K$ and $d$. Therefore
\[
    \overline K(s)
    =\E\bigl[(K-\abs d)\1_{\{\abs d\le K\}}\bigr]
    =K-\E\abs d+O(e^{-cK^2})=K+O(1).
\]
Also, by \eqref{eq:xi-mean},
\[
    \E d
    =
    \E\xi-\frac12
    =
    -\frac12+O(K^{-2}),
\]
and \eqref{eq:Fp-local}, together with the exponentially small tail
contribution (here $F'\leq K^2$ on the complement and
$\E[K^2\1_{\{\abs d>K\}}]=O(1)$), gives
\[
    \E F'(\rho+sT)
    =
    \frac{K^2}{2}
    +
    K\,\E d
    -
    \frac12\,\E[d\abs d]
    +
    O(1)
    =
    \frac{K^2}{2}
    +
    K\left(-\frac12+O(K^{-2})\right)
    +
    O(1)
    =
    \rho_0-\frac K2+O(1),
\]
using $\abs{\E[d\abs d]}\leq\E d^2=O(1)$.

It remains to estimate $G(s)$. Decompose \eqref{eq:G-def} as
$G(s)=(\mathrm I)+(\mathrm{II})-(\mathrm{III})$ with
\[
    (\mathrm I)=\frac1m\E\!\left[F''\frac{r^2}{\rho^2}\right],
    \quad
    (\mathrm{II})=\E\!\left[\frac{F'}{\rho}\right],
    \quad
    (\mathrm{III})=\frac1m\E\!\left[F'\frac{r^2}{\rho^3}\right].
\]
The first term satisfies $0\le(\mathrm I)\leq K/m=O(K^{-3})$, since
$F''\leq K$ and $r^2/\rho^2\leq1$. For the second, using
\[
    \left\|
        \frac1\rho-\frac1{\rho_0}
    \right\|_{L^2}
    =
    \left\|
        \frac{\xi}{\rho\rho_0}
    \right\|_{L^2}
    \le\frac{\norm\xi_{L^2}}{\sqrt2\,m}
    =
    O(K^{-4})
\]
(because $\rho\ge\sqrt m$, $\rho_0=\sqrt{2m}$, and $\norm\xi_{L^2}=O(1)$) and
$F'\leq K^2$,
\[
    (\mathrm{II})
    =
    \E\!\left[\frac{F'}{\rho_0}\right]
    +
    K^2\cdot O(K^{-4})
    =
    \frac{\E F'}{\rho_0}
    +
    O(K^{-2}).
\]
Finally $r\mapsto r^2/(m+r^2)^{3/2}$ is maximized at $r^2=2m$ with value
$2\cdot3^{-3/2}m^{-1/2}\leq m^{-1/2}$, so
$0\le(\mathrm{III})\leq K^2m^{-3/2}=O(K^{-4})$. Combining with
\eqref{eq:Fprime-est} and $\rho_0=K^2/2$,
\[
    G(s)
    =
    \frac{\E F'(\rho+sT)}{\rho_0}
    +
    O(K^{-2})
    =
    1-\frac{K/2}{\rho_0}+O(\rho_0^{-1})+O(K^{-2})
    =
    1-\frac1K+O(K^{-2}).
    \qedhere
\]
\end{proof}

\begin{lemma}[Exact isotropic tuning]
\label{lem:tuning}
For all sufficiently large $K$, there exists
\[
    s=s_K\in(0,2K^{-1/2})
\]
such that
\begin{equation}
    s^2\overline K(s)=G(s).
    \label{eq:tuning}
\end{equation}
Writing
\begin{equation}
    a:=s^2\overline K(s)=G(s),
    \qquad
    \alpha:=1-a,
    \label{eq:a-alpha}
\end{equation}
one has
\begin{align}
    s^2
    &=
    K^{-1}+O(K^{-2}),
    \label{eq:s-est}\\
    a
    &=
    1-K^{-1}+O(K^{-2}),
    \label{eq:a-est}\\
    \alpha
    &=
    K^{-1}+O(K^{-2}),
    \label{eq:alpha-est}
\end{align}
and
\begin{equation}
    \E A_s=a\Id_n.
    \label{eq:exact-mean-As}
\end{equation}
\end{lemma}

\begin{proof}
The function
\[
    \Phi_K(s):=s^2\overline K(s)-G(s)
\]
is continuous on $[0,2K^{-1/2}]$. Indeed, the integrands are pointwise
continuous in $s$. The integrand defining $\overline K(s)$ is bounded by
$K$, while the absolute value of the full integrand defining $G(s)$ is
bounded by
\[
    \frac Km
    +
    K^2\left(\frac1\rho+\frac{r^2}{m\rho^3}\right)
    \leq
    \frac Km+\frac{2K^2}{\sqrt m}.
\]
Thus dominated convergence applies. Now $\Phi_K(0)=-G(0)$,
and $G(0)=1-K^{-1}+O(K^{-2})>0$ for large $K$ by \eqref{eq:G-est}, so
$\Phi_K(0)<0$. At $s=2K^{-1/2}$, \Cref{lem:profile-est} gives
\[
    s^2\overline K(s)=\frac4K\bigl(K+O(1)\bigr)=4+O(K^{-1}),
    \qquad
    G(s)=1+O(K^{-1}),
\]
so $\Phi_K(2K^{-1/2})=3+O(K^{-1})>0$. The intermediate value theorem yields a
zero $s_K\in(0,2K^{-1/2})$.

At this zero, \eqref{eq:G-est} gives $a=G(s_K)=1-K^{-1}+O(K^{-2})$, hence
\eqref{eq:a-est} and \eqref{eq:alpha-est}; in particular $0<\alpha<1$ for
large $K$. Dividing $s_K^2\overline K(s_K)=a$ by \eqref{eq:barK-est} yields
$s_K^2=a/\overline K(s_K)=(1+O(K^{-1}))/K$, which is \eqref{eq:s-est}.
Finally, \eqref{eq:mean-As} and the tuning equation
$s_K^2\overline K(s_K)=G(s_K)=a$ give \eqref{eq:exact-mean-As}.
\end{proof}

\subsection{The potential}

Fix the tuned value $s=s_K$ from \Cref{lem:tuning}. Since
$\abs{\nabla U_s}=F'(x)\norm g\leq\sqrt2\,K^2$ is bounded, the vector
\begin{equation}
    \ell:=\E[\nabla U_s(Z)]
    \label{eq:ell-def}
\end{equation}
is well defined, and we set
\begin{equation}
    V(z):=\frac{\alpha}{2}\norm z^2+U_s(z)-\ell^\top z.
    \label{eq:V-def}
\end{equation}

\begin{lemma}[Exact whitened equilibrium]
\label{lem:whitening}
The potential \eqref{eq:V-def} satisfies
\[
    \E[\nabla V(Z)]=0,
    \qquad
    \E[\nabla^2V(Z)]=\Id_n.
\]
\end{lemma}

\begin{proof}
Since $\E Z=0$,
\[
    \E[\nabla V(Z)]
    =
    \alpha\E Z+\E[\nabla U_s(Z)]-\ell
    =
    0.
\]
Also, by \eqref{eq:exact-mean-As},
\[
    \E[\nabla^2V(Z)]
    =
    \alpha\Id_n+\E A_s
    =
    (\alpha+a)\Id_n
    =
    \Id_n.
\]
\end{proof}

% ============================================================
\section{Ellipticity and a dimension-free Lipschitz Hessian}
% ============================================================

\begin{lemma}[Uniform ellipticity and condition number]
\label{lem:ellipticity}
Let
\[
    L_K:=\sup_{z\in\R^n}\norm{A_s(z)}_{\op},
    \qquad
    \beta:=\alpha+L_K.
\]
For all sufficiently large $K$,
\begin{equation}
    \frac K2
    \leq
    L_K
    \leq
    3K.
    \label{eq:LK-bounds}
\end{equation}
Consequently,
\[
    \alpha\Id_n
    \preceq
    \nabla^2V(z)
    \preceq
    \beta\Id_n,
    \qquad z\in\R^n,
\]
and
\begin{equation}
    \kappa=\frac{\beta}{\alpha}=\Theta(K^2).
    \label{eq:kappa-K}
\end{equation}
\end{lemma}

\begin{proof}
Since $\rho\geq\sqrt m=K^2/\sqrt8$ and
$s^2\leq4/K$, for large $K$,
\[
    \norm g^2
    =
    s^2+\frac{r^2}{\rho^2}
    \leq
    1+s^2
    \leq2,
    \qquad
    \norm{D^2\rho}_{\op}
    \leq
    \frac1\rho
    \leq
    \frac{\sqrt8}{K^2}.
\]
Using \eqref{eq:F-basic} in \eqref{eq:As},
\[
    0\preceq A_s
    \preceq
    \left(F''\norm g^2+F'\norm{D^2\rho}_{\op}\right)\Id_n
    \preceq
    \left(2K+\sqrt8\right)\Id_n
    \preceq
    3K\Id_n
\]
for large $K$. This proves the upper bound in \eqref{eq:LK-bounds}.

For the lower bound, choose $y$ with $r^2=m$, so $\rho=\rho_0$, and
choose $t=(c-\rho_0)/s=1/(2s)$, so that $x=c$. Then $F''(x)=K$ and
\[
    \norm g^2
    =
    s^2+\frac{r^2}{\rho_0^2}
    =
    s^2+\frac12
    \geq
    \frac12.
\]
Therefore the rank-one matrix $F''(x)gg^\top=Kgg^\top$ has an eigenvalue
$K\norm g^2\geq K/2$, and since the second summand in \eqref{eq:As} is
positive semidefinite, $\norm{A_s(z)}_{\op}\geq K/2$ at this point; hence
$L_K\geq K/2$.

Since $A_s\succeq0$, the lower Hessian bound is $\alpha\Id_n$. It is
also attained: for any fixed $y$, choose $t$ sufficiently negative so
that $\rho(y)+st\leq c-K$; then $F'=F''=0$ and $A_s=0$, so no smaller
$\alpha$ is admissible. Finally, \eqref{eq:alpha-est} and
\eqref{eq:LK-bounds} give $\kappa=\beta/\alpha=(\alpha+L_K)/\alpha=\Theta(K^2)$,
which is \eqref{eq:kappa-K}.
\end{proof}

For a trilinear form $B$, write
\[
    \norm B_{\inj}
    :=
    \sup_{\norm u=\norm v=\norm w=1}\abs{B[u,v,w]}.
\]

\begin{lemma}[Dimension-free Hessian Lipschitz bound]
\label{lem:H3}
The potential \eqref{eq:V-def} belongs to $C^{2,1}(\R^n)$ and satisfies
\begin{equation}
    \norm{\nabla^2V(x)-\nabla^2V(y)}_{\op}
    \leq
    64\norm{x-y},
    \qquad x,y\in\R^n.
    \label{eq:H3-64}
\end{equation}
\end{lemma}

\begin{proof}
The function $F''$ is globally Lipschitz and differentiable almost
everywhere, with $\abs{F'''}\leq1$. At every point of differentiability,
and for unit vectors $u,v,w\in\R^n$,
\begin{align}
    D^3U_s[u,v,w]
    ={}&
    F'''(x)g[u]g[v]g[w]
    \nonumber\\
    &+
    F''(x)\Bigl(
        g[u]D^2\rho[v,w]
        +
        g[v]D^2\rho[u,w]
        +
        g[w]D^2\rho[u,v]
    \Bigr)
    \nonumber\\
    &+
    F'(x)D^3\rho[u,v,w].
    \label{eq:D3U}
\end{align}
Here $\rho$ is viewed as a function of $(t,y)$ independent of $t$, and the
middle line collects the three product-rule terms in which exactly one
differentiation falls on the factor $gg^\top$.

A direct differentiation gives, for $u,v,w\in\R^m$,
\begin{align}
    D^3\rho[u,v,w]
    ={}&
    -\frac{
        (u^\top v)(y^\top w)
        +(u^\top w)(y^\top v)
        +(v^\top w)(y^\top u)
    }{\rho^3}
    \nonumber\\
    &+
    \frac{
        3(y^\top u)(y^\top v)(y^\top w)
    }{\rho^5}.
    \label{eq:D3rho}
\end{align}
Therefore
\[
    \norm{D^3\rho}_{\inj}
    \leq
    \frac{3r}{\rho^3}+\frac{3r^3}{\rho^5}
    \leq
    \frac6{\rho^2}
    \leq
    \frac6m
    =
    \frac{48}{K^4}.
\]
Using $\norm g\leq\sqrt2$,
$\norm{D^2\rho}_{\op}\leq\sqrt8/K^2$, and
\eqref{eq:F-basic}, \eqref{eq:D3U} gives, almost everywhere,
\[
    \norm{D^3U_s}_{\inj}
    \leq
    \abs{F'''}\norm g^3
    +3F''\norm g\,\norm{D^2\rho}_{\op}
    +F'\norm{D^3\rho}_{\inj}
    \leq
    2\sqrt2+\frac{12}{K}+\frac{48}{K^2}
    \leq
    64
\]
for all sufficiently large $K$.

We now integrate this almost-everywhere bound into a global Lipschitz bound.
Fix unit vectors $u,v\in\R^n$ and set
$f_{u,v}(z):=u^\top A_s(z)v$. Each entry of $A_s=D^2U_s$ is locally Lipschitz
(the profile $F''$ is Lipschitz, $F'$ is $C^{1,1}$, and $\rho,g,D^2\rho$ are
smooth on $\{\rho\geq\sqrt m\}$), so $f_{u,v}$ is locally Lipschitz, hence
differentiable almost everywhere with
$\abs{\nabla f_{u,v}(z)}=\abs{D^3U_s(z)[u,v,\cdot]}\leq\norm{D^3U_s(z)}_{\inj}
\leq64$ wherever it exists. For a locally Lipschitz function on $\R^n$ the
global Lipschitz constant equals the essential supremum of $\abs{\nabla f}$:
mollifying, $f_{u,v}*\eta_\varepsilon$ is smooth with
$\abs{\nabla(f_{u,v}*\eta_\varepsilon)}\leq64$ everywhere, hence
$64$-Lipschitz, and letting $\varepsilon\downarrow0$ preserves the
inequality. Thus
\[
    \abs{u^\top\bigl(A_s(x)-A_s(y)\bigr)v}\leq64\norm{x-y},
\]
and taking the supremum over unit $u,v$ yields
$\norm{A_s(x)-A_s(y)}_{\op}\leq64\norm{x-y}$. The quadratic and linear terms
in \eqref{eq:V-def} contribute a constant Hessian $\alpha\Id_n$, which does
not affect the Hessian Lipschitz seminorm, proving \eqref{eq:H3-64}.
\end{proof}

% ============================================================
\section{A large first-Hermite coupling}
% ============================================================

Define
\begin{equation}
    X
    :=
    \diag\left(0,\frac{\Id_m}{\sqrt m}\right)
    \in\Sym(n),
    \qquad
    \norm X_{\mathrm F}=1.
    \label{eq:X-def}
\end{equation}

\begin{lemma}[Codazzi conversion of the collective radial mode]
\label{lem:codazzi}
The vector $T[X]$ points in the distinguished $t$ direction and
\begin{equation}
    \norm{T[X]}
    =
    \frac{s}{\sqrt m}
    \E\!\left[
        F''(\rho+sT)\frac{r^2}{\rho}
    \right].
    \label{eq:TX-exact}
\end{equation}
\end{lemma}

\begin{proof}
The constant matrix $\alpha\Id_n$ has zero first-Hermite coefficient
($\E[\Tr(X\cdot\alpha\Id_n)Z]=\alpha\sqrt m\,\E Z=0$, using
$\Tr X=\sqrt m$), so only $A_s=D^2U_s$ contributes:
$T[X]=\E[\Tr(XA_s)Z]$. Now
\[
    \Tr(XA_s)
    =\frac1{\sqrt m}\sum_{i=1}^m\partial_{ii}U_s
    =\frac1{\sqrt m}\left(
        F''\frac{r^2}{\rho^2}+F'\Bigl(\frac m\rho-\frac{r^2}{\rho^3}\Bigr)
    \right)
\]
is a function of $(T,r)$ alone. For each $i$, the reflection
$y_i\mapsto-y_i$ preserves the law of $Z$ and the value of $\Tr(XA_s)$ while
flipping $Y_i$, so $\E[\Tr(XA_s)Y_i]=0$; hence every $y$-component of $T[X]$
vanishes, and $T[X]=\E[\Tr(XA_s)T]\,e_t$. Its $t$-component is
\[
    \frac1{\sqrt m}
    \sum_{i=1}^m
    \E\!\left[
        T\,\partial_{ii}U_s(T,Y)
    \right].
\]
Let $\varphi_n$ denote the standard Gaussian density. Since $A_s=D^2U_s$
is bounded, its entries define tempered distributions. Commutation of
distributional derivatives gives
\[
    \partial_t(A_s)_{ii}=\partial_i(A_s)_{ti}.
\]
Consequently,
\begin{align*}
    \E[T(A_s)_{ii}(T,Y)]
    &=
    -\int_{\R^n}(A_s)_{ii}\,\partial_t\varphi_n\\
    &=
    \left\langle\partial_t(A_s)_{ii},\varphi_n\right\rangle\\
    &=
    \left\langle\partial_i(A_s)_{ti},\varphi_n\right\rangle\\
    &=
    -\int_{\R^n}(A_s)_{ti}\,\partial_i\varphi_n\\
    &=
    \E[Y_i(A_s)_{ti}(T,Y)].
\end{align*}
These pairings may equivalently be justified by first inserting compactly
supported cutoffs of $\varphi_n$ and then letting the cutoff radius tend to
infinity. From \eqref{eq:As},
\[
    (A_s)_{ti}
    =
    sF''(\rho+sT)\frac{Y_i}{\rho}.
\]
Summing over $i$ therefore yields
\[
    e_t^\top T[X]
    =
    \frac{s}{\sqrt m}
    \E\!\left[F''(\rho+sT)\frac{r^2}{\rho}\right].
\]
The right-hand side is nonnegative, proving \eqref{eq:TX-exact}.
\end{proof}

\begin{lemma}[Size of the coupling]
\label{lem:TX-size}
There is a numerical constant $C$ such that
\begin{equation}
    \norm{T[X]}^2
    =
    \frac K2+O(1),
    \qquad
    \tau_H^2
    \geq
    \frac K2-C.
    \label{eq:tau-lower}
\end{equation}
\end{lemma}

\begin{proof}
Since $r^2=\rho^2-m$, define
\[
    q(\rho):=\frac{r^2}{\rho}=\rho-\frac m\rho.
\]
On $[\sqrt m,\infty)$,
\[
    q'(\rho)=1+\frac m{\rho^2}\leq2,
\]
so $q$ is $2$-Lipschitz there, and
\[
    q(\rho_0)
    =
    \rho_0-\frac m{\rho_0}
    =
    \frac{\rho_0}{2}
    =
    \frac{K^2}{4}.
\]
Therefore, using $0\leq F''\leq K$, $\abs{q(\rho)-q(\rho_0)}\leq2\abs\xi$,
and \Cref{lem:xi},
\begin{align}
    \E[F''q(\rho)]
    &=
    q(\rho_0)\,\E F''
    +
    O\!\left(K\,\E\abs{\rho-\rho_0}\right)
    \nonumber\\
    &=
    \overline K(s)\frac{K^2}{4}
    +
    O(K).
    \label{eq:Fppq}
\end{align}
Because
\[
    \sqrt m=\frac{K^2}{2\sqrt2},
    \qquad\text{so}\qquad
    \frac1m\cdot\frac{K^4}{16}=\frac12,
\]
substitution into \eqref{eq:TX-exact} yields
\begin{align*}
    \norm{T[X]}^2
    &=
    \frac{s^2}{m}
    \left(
        \overline K(s)\frac{K^2}{4}+O(K)
    \right)^2\\
    &=
    \frac{s^2\overline K(s)^2K^4}{16m}
    +O\!\left(\frac{s^2\overline K(s)K^3}{m}\right)
    +O\!\left(\frac{s^2K^2}{m}\right)\\
    &=
    \frac{s^2\overline K(s)^2}{2}+O(K^{-1})\\
    &=
    \frac{a\,\overline K(s)}{2}+O(K^{-1})\\
    &=
    \frac K2+O(1),
\end{align*}
where we used $s^2\overline K(s)=a$ (\Cref{lem:tuning}),
$\overline K(s)=K+O(1)$ and $a=1+O(K^{-1})$ (\Cref{lem:profile-est,lem:tuning}),
and $s^2=O(K^{-1})$.

Finally, $\norm X_{\mathrm F}=1$ and
$\tau_H=\norm T_{\Sym(n)\to\R^n}$ by \eqref{eq:tau-def}, so
\[
    \tau_H^2
    \geq
    \norm{T[X]}^2
    \geq
    \frac K2-C.
    \qedhere
\]
\end{proof}

\begin{proof}[Proof that the logarithmic coupling estimate fails]
By \Cref{lem:tuning}, $a=1+O(K^{-1})\in[\tfrac12,1]$ for large $K$, and by
\Cref{lem:ellipticity}, $L=L_K\leq3K$. Hence
\[
    a^2\left(1+\log\frac{eL}{a}\right)
    \leq
    1+\log(6eK)
    =
    O(\log K).
\]
On the other hand, \eqref{eq:tau-lower} gives
$\tau_H^2\geq K/2-C$. Therefore no absolute constant can make
\[
    \tau_H^2
    \leq
    C a^2\left(1+\log\frac{eL}{a}\right)
\]
valid for all dimensions.
\end{proof}

\begin{remark}[Direct cross-check of \eqref{eq:TX-exact}]
\label{rem:crosscheck}
The identity can also be checked directly, without commuting mixed weak
derivatives. Gaussian integration by parts in $T$ gives
\[
    e_t^\top T[X]
    =\frac{s}{\sqrt m}\,
    \E\!\left[F'''\frac{r^2}{\rho^2}
        +F''\Bigl(\frac m\rho-\frac{r^2}{\rho^3}\Bigr)\right].
\]
Here $F'''$ denotes the almost-everywhere derivative of the Lipschitz
function $F''$. For fixed $T$, apply Gaussian integration by parts in $Y$
to the vector field
\[
    \Psi_T(y)
    :=
    F''(\rho(y)+sT)\frac{y}{\rho(y)}.
\]
Since
\[
    Y^\top\Psi_T(Y)=F''\frac{r^2}{\rho}
\]
and
\[
    \operatorname{div}\Psi_T
    =
    F'''\frac{r^2}{\rho^2}
    +
    F''\left(\frac m\rho-\frac{r^2}{\rho^3}\right),
\]
the Gaussian divergence identity
$\E[Y^\top\Psi_T(Y)]=\E[\operatorname{div}\Psi_T(Y)]$ recovers
\eqref{eq:TX-exact} exactly.

This direct calculation verifies the identity; it does not distinguish
Hessian fields from non-gradient matrix fields. Indeed, $T[X]$ depends only
on $\Tr(XW)$, and for the present $X$ any matrix field with the same $yy$
block has the same value of $T[X]$, independently of its $ty$ block. Finally,
\Cref{lem:TX-size,lem:ellipticity} imply
\[
    \norm{T[X]}^2=\Theta(K)=\Theta(\kappa^{1/2}),
    \qquad
    \norm{T[X]}=\Theta(K^{1/2})=\Theta(\kappa^{1/4}).
\]
\end{remark}

% ============================================================
\section{A nearly-null direction for the packed local generator}
\label{sec:gap}
% ============================================================

We first record the exact second-Hermite identity, in polarized form, from
which self-adjointness of $H$ also follows.

\begin{lemma}[Second-Hermite/Bochner identity]
\label{lem:second-Hermite}
Let $V\in C^2(\R^n)$, assume that $W=\nabla^2V$ is bounded, and suppose
$\E W(Z)=\Id_n$. Then for all $X,Y\in\Sym(n)$,
\begin{equation}
    \inner{Y}{H[X]}_{\mathrm F}
    =
    \E\!\left[
        (YZ)^\top\bigl(W(Z)-\Id_n\bigr)(XZ)
    \right].
    \label{eq:second-Hermite-id}
\end{equation}
In particular the right-hand side is symmetric in $(X,Y)$ (because $W-\Id_n$
is a symmetric matrix), so $H$ is self-adjoint on $\Sym(n)$; and, taking
$Y=X$,
\begin{equation}
    \inner{X}{H[X]}_{\mathrm F}
    =
    \E\!\left[
        (XZ)^\top\bigl(W(Z)-\Id_n\bigr)(XZ)
    \right].
    \label{eq:second-Hermite-quad}
\end{equation}
\end{lemma}

\begin{proof}
Let $\varphi_n$ denote the standard Gaussian density. Since $W$ is bounded,
its entries define tempered distributions, and
\[
    \partial_{cd}\varphi_n(z)
    =
    (z_cz_d-\delta_{cd})\varphi_n(z).
\]
Because $W=D^2V$, commutation of distributional derivatives gives
\[
    \partial_{cd}W_{ab}
    =
    \partial_{bc}W_{ad}.
\]
Therefore
\begin{align*}
    \E[(Z_cZ_d-\delta_{cd})W_{ab}(Z)]
    &=
    \left\langle W_{ab},\partial_{cd}\varphi_n\right\rangle\\
    &=
    \left\langle\partial_{cd}W_{ab},\varphi_n\right\rangle\\
    &=
    \left\langle\partial_{bc}W_{ad},\varphi_n\right\rangle\\
    &=
    \left\langle W_{ad},\partial_{bc}\varphi_n\right\rangle\\
    &=
    \E[(Z_bZ_c-\delta_{bc})W_{ad}(Z)].
\end{align*}
One may instead insert compactly supported cutoffs of $\varphi_n$; boundedness
of $W$ makes all cutoff-error terms vanish in the limit.

By \eqref{eq:H-def} and $\E[Z_cZ_d-\delta_{cd}]=0$ (which annihilates the
constant matrix $-\Id_n$),
\begin{align}
    \inner{Y}{H[X]}_{\mathrm F}
    &=
    \sum_{a,b,c,d}
    Y_{ab}X_{cd}
    \E\!\left[
        (Z_cZ_d-\delta_{cd})\bigl(W_{ab}(Z)-\delta_{ab}\bigr)
    \right]
    \nonumber\\
    &=
    \sum_{a,b,c,d}
    Y_{ab}X_{cd}
    \E\!\left[
        (Z_cZ_d-\delta_{cd})W_{ab}(Z)
    \right].
    \label{eq:h-indices-1}
\end{align}
Substituting into \eqref{eq:h-indices-1},
\begin{align*}
    \inner{Y}{H[X]}_{\mathrm F}
    &=
    \E\!\left[
        \sum_{a,b,c,d}
        Y_{ab}X_{cd}Z_bZ_cW_{ad}
    \right]
    -
    \sum_{a,b,c,d}
    Y_{ab}X_{cd}\delta_{bc}\,\E W_{ad}\\
    &=
    \E[(YZ)^\top W(Z)(XZ)]
    -
    \Tr(YX),
\end{align*}
where we used $(YZ)_a=\sum_bY_{ab}Z_b$, $(XZ)_d=\sum_cX_{cd}Z_c$, and
$\sum_bY_{ab}X_{bd}=(YX)_{ad}$, together with $\E W(Z)=\Id_n$. Since
$\Tr(YX)=\E[(YZ)^\top(XZ)]$, this proves
\eqref{eq:second-Hermite-id}; the symmetry and
\eqref{eq:second-Hermite-quad} are immediate.
\end{proof}

For the normalized matrix $X$ in \eqref{eq:X-def}, define
\begin{equation}
    \tau_X^2:=\norm{T[X]}^2,
    \qquad
    h_X:=\inner{X}{H[X]}_{\mathrm F}.
    \label{eq:tau-h-def}
\end{equation}

\begin{lemma}[Second-Hermite asymptotics]
\label{lem:hX}
For the matrix $X$ in \eqref{eq:X-def},
\begin{equation}
    h_X
    =
    \frac K2+O(1),
    \qquad
    h_X-\tau_X^2
    =
    O(1).
    \label{eq:hX-est}
\end{equation}
\end{lemma}

\begin{proof}
Since
\[
    XZ=\left(0,\frac{Y}{\sqrt m}\right)
    \qquad\text{and}\qquad
    W-\Id_n
    =
    \alpha\Id_n+A_s-\Id_n
    =
    A_s-a\Id_n,
\]
\eqref{eq:second-Hermite-quad} gives
$h_X=\frac1m\E[Y^\top(A_s)_{yy}Y]-a$, where $(A_s)_{yy}$ is the $yy$ block of
\eqref{eq:As}:
\[
    (A_s)_{yy}=F''\frac{yy^\top}{\rho^2}
    +F'\left(\frac{\Id_m}{\rho}-\frac{yy^\top}{\rho^3}\right).
\]
Using $Y^\top yy^\top Y=r^4$ (with $y=Y$), $Y^\top Y=r^2$, and
$\rho^2-r^2=m$,
\[
    Y^\top(A_s)_{yy}Y
    =F''\frac{r^4}{\rho^2}
    +F'\left(\frac{r^2}{\rho}-\frac{r^4}{\rho^3}\right)
    =F''\frac{r^4}{\rho^2}+F'\frac{m\,r^2}{\rho^3},
\]
so that
\begin{equation}
    h_X
    =
    \frac1m
    \E\!\left[
        F''\frac{r^4}{\rho^2}
    \right]
    +
    \E\!\left[
        F'\frac{r^2}{\rho^3}
    \right]
    -
    a.
    \label{eq:hX-exact}
\end{equation}

Let $u:=r^2/m$. Then
\[
    \frac{r^4}{m\rho^2}
    =
    \frac{(mu)^2}{m(m+mu)}
    =
    \frac{u^2}{1+u}.
\]
The function $\phi(u)=u^2/(1+u)$ satisfies $0\le\phi'(u)=1-(1+u)^{-2}<1$ on
$[0,\infty)$, hence is $1$-Lipschitz with $\phi(1)=1/2$. Moreover
\[
    \norm{u-1}_{L^2}
    =
    \frac{\norm{r^2-m}_{L^2}}{m}
    =
    \frac{\sqrt{2m}}{m}
    =
    \sqrt{\frac2m}
    =
    O(K^{-2}).
\]
Therefore
\begin{align*}
    \frac1m
    \E\!\left[
        F''\frac{r^4}{\rho^2}
    \right]
    &=
    \E[F''\phi(u)]
    =
    \frac12\E F''
    +
    \E\bigl[F''(\phi(u)-\tfrac12)\bigr]\\
    &=
    \frac{\overline K(s)}2
    +
    O\!\left(K\,\E\abs{u-1}\right)
    =
    \frac{\overline K(s)}2+O(K^{-1})
    =
    \frac K2+O(1).
\end{align*}
For the second term in \eqref{eq:hX-exact}, since
$\sup_{r\geq0}r^2(m+r^2)^{-3/2}\leq m^{-1/2}$ and $F'\leq K^2$,
\[
    0
    \leq
    \E\!\left[
        F'\frac{r^2}{\rho^3}
    \right]
    \leq
    \frac{K^2}{\sqrt m}
    =
    2\sqrt2
    =
    O(1).
\]
Since $a=O(1)$, this proves $h_X=K/2+O(1)$.

By \Cref{lem:TX-size} and $s^2\overline K=a$,
$\tau_X^2=\frac{s^2\overline K^2}2+O(K^{-1})=\frac{a\overline K}2+O(K^{-1})
=\frac{aK}2+O(1)$. Hence
\[
    h_X-\tau_X^2
    =
    \frac{(1-a)\overline K}{2}+O(1)
    =
    \frac{\alpha K}{2}+O(1)
    =
    O(1),
\]
using \eqref{eq:alpha-est}.
\end{proof}

\begin{lemma}[Explicit Rayleigh direction]
\label{lem:Rayleigh}
Set
\begin{equation}
    Y:=\frac{X}{\sqrt2},
    \qquad
    v:=-\frac{T[X]}{2}.
    \label{eq:Rayleigh-vector}
\end{equation}
Then
\begin{equation}
    \frac{
        \inner{(v,Y)}{L_\star(v,Y)}_{\mathrm{pack}}
    }{
        \norm v^2+\norm Y_{\mathrm F}^2
    }
    =
    \frac{2+h_X-\tau_X^2}{2+\tau_X^2}.
    \label{eq:Rayleigh-quotient}
\end{equation}
\end{lemma}

\begin{proof}
Since $\norm X_{\mathrm F}=1$,
\[
    \norm v^2=\frac{\tau_X^2}{4},
    \qquad
    \norm Y_{\mathrm F}^2=\frac12.
\]
By \eqref{eq:Lstar-def}, the packed quadratic form is
\[
    \inner{(v,Y)}{L_\star(v,Y)}_{\mathrm{pack}}
    =
    \norm v^2
    +
    2\inner{v}{(T/\sqrt2)Y}
    +
    \norm Y_{\mathrm F}^2
    +
    \frac12\inner{Y}{H[Y]}_{\mathrm F},
\]
where the two off-diagonal contributions combine into
$2\inner{v}{(T/\sqrt2)Y}$ by self-adjointness of $L_\star$. Now
$(T/\sqrt2)Y=T[X]/2$ and, by \eqref{eq:tau-h-def},
$\inner{Y}{H[Y]}_{\mathrm F}=\tfrac12\inner{X}{H[X]}_{\mathrm F}=h_X/2$.
Hence
\[
    \inner{(v,Y)}{L_\star(v,Y)}_{\mathrm{pack}}
    =
    \frac{\tau_X^2}{4}
    +
    2\cdot\Bigl(-\frac{T[X]}2\Bigr)^{\!\top}\frac{T[X]}2
    +
    \frac12
    +
    \frac{h_X}{4}
    =
    \frac{\tau_X^2}{4}
    -
    \frac{\tau_X^2}{2}
    +
    \frac12
    +
    \frac{h_X}{4}
    =
    \frac{2+h_X-\tau_X^2}{4}.
\]
The denominator is
\[
    \norm v^2+\norm Y_{\mathrm F}^2
    =
    \frac{\tau_X^2}{4}+\frac12
    =
    \frac{2+\tau_X^2}{4},
\]
which proves \eqref{eq:Rayleigh-quotient}.
\end{proof}

\begin{proof}[Proof of the local-gap upper bound]
Since $L_\star$ is self-adjoint, the Rayleigh--Ritz variational principle and
\Cref{lem:Rayleigh} give
\[
    \gamma_{\mathrm{loc}}
    =
    \lambda_{\min}(L_\star)
    \leq
    \frac{2+h_X-\tau_X^2}{2+\tau_X^2}.
\]
By \Cref{lem:hX}, the numerator is bounded above by a numerical constant
$2+C_1$. By \Cref{lem:TX-size}, $\tau_X^2\geq K/2-C$, so the denominator is
at least a positive multiple of $K$ for large $K$. Hence
\[
    \gamma_{\mathrm{loc}}
    \leq
    \frac{C_1}{K}.
\]
Together with $\kappa=\Theta(K^2)$ from \Cref{lem:ellipticity}, this gives
$\gamma_{\mathrm{loc}}\leq C_1'/\sqrt\kappa$.

Every member of the sequence satisfies the Hessian-Lipschitz condition
\eqref{eq:H3} with the common admissible constant $L_H=64$ by
\Cref{lem:H3}. If a positive function $c(L_H)$ satisfied
\[
    \gamma_{\mathrm{loc}}
    \geq
    \frac{c(L_H)}{1+\log\kappa}
\]
in every dimension, then applying it with $L_H=64$ would give
\[
    \frac{c(64)}{1+\log\kappa}
    \leq
    \frac{C_1}{K}.
\]
But $\log\kappa=\Theta(\log K)$ and
$K^{-1}=o((\log K)^{-1})$, a contradiction for large $K$.
\end{proof}

\begin{proof}[Proof of \Cref{thm:main}]
Fix $m$ with $K=(8m)^{1/4}$ large and let $V$ be the potential
\eqref{eq:V-def}. \Cref{lem:whitening} gives the whitening conditions
\eqref{eq:H2}, and \Cref{lem:ellipticity} gives \eqref{eq:H1} together with
\eqref{eq:kappa-main} (with $c_0,C_0$ from $\kappa=\Theta(K^2)$).
\Cref{lem:TX-size} gives \eqref{eq:tau-main}, and its corollary above shows
the failure of the logarithmic coupling estimate, i.e.\ conclusion~(1).
\Cref{lem:H3} gives \eqref{eq:H3} with $L_H=64$. The local-gap upper bound
above gives \eqref{eq:gap-main}, and hence conclusion~(2). This proves the
theorem.
\end{proof}

% ============================================================
\section{Scope, regularity, and smooth realization}
% ============================================================

\begin{corollary}[Smooth realization]
\label{cor:smooth-ridge}
After increasing the numerical threshold $K_0$ if necessary, the potentials
in \Cref{thm:main} may be chosen in $C^\infty(\R^n)$ while retaining all
conclusions of the theorem and the common Hessian-Lipschitz bound $L_H=64$.
\end{corollary}

\begin{proof}
Let $f_K(x):=(K-\abs{x-c})_+$ and fix an even, nonnegative
$\eta\in C_c^\infty((-1,1))$ with $\int\eta=1$. For a fixed
$\varepsilon\in(0,1]$, set
\[
    \eta_\varepsilon(u):=\varepsilon^{-1}\eta(u/\varepsilon),
    \qquad
    f_{K,\varepsilon}:=f_K*\eta_\varepsilon,
\]
and define $F'_{K,\varepsilon}(-\infty)=0$ with
$F''_{K,\varepsilon}=f_{K,\varepsilon}$. Then
\[
    0\leq f_{K,\varepsilon}\leq K,
    \qquad
    \Lip(f_{K,\varepsilon})\leq1,
    \qquad
    \int_\R f_{K,\varepsilon}=K^2,
\]
and
\[
    \norm{f_{K,\varepsilon}-f_K}_{L^\infty}\leq\varepsilon.
\]
Moreover, $F'_{K,\varepsilon}=F'_K*\eta_\varepsilon$. Since $F'_K$ has a
$1$-Lipschitz derivative and $\eta_\varepsilon$ has zero first moment,
Taylor's formula with Lipschitz remainder gives
\[
    \norm{F'_{K,\varepsilon}-F'_K}_{L^\infty}
    \leq
    \frac{\varepsilon^2}{2}.
\]
Consequently, uniformly for $0\leq s\leq2K^{-1/2}$,
\[
    \abs{\overline K_\varepsilon(s)-\overline K(s)}\leq\varepsilon,
    \qquad
    \abs{\E F'_{K,\varepsilon}(\rho+sT)-\E F'_K(\rho+sT)}
    \leq
    \frac{\varepsilon^2}{2},
\]
and
\[
    \abs{G_\varepsilon(s)-G(s)}
    \leq
    \frac{\varepsilon}{m}
    +
    \frac{\varepsilon^2}{2}
    \left(m^{-1/2}+m^{-3/2}\right).
\]
Thus all estimates in \Cref{lem:profile-est} remain unchanged. The same
intermediate-value argument retunes $s=s_{K,\varepsilon}$ and restores exact
whitening. Also $f_{K,\varepsilon}(c)\geq K-\varepsilon$, so the ellipticity
lower bound remains of order $K$. Finally,
$\norm{f'_{K,\varepsilon}}_{L^\infty}\leq1$, and the proof of
\Cref{lem:H3} gives the same admissible constant $L_H=64$. Repeating the
coupling and Rayleigh calculations proves all asserted asymptotics.
\end{proof}

\begin{remark}[Scaling and dimension scope]
The construction satisfies
\[
    s^2=\Theta(K^{-1}),
    \qquad
    \alpha=\Theta(K^{-1}),
    \qquad
    \beta=\Theta(K),
    \qquad
    \kappa=\Theta(K^2),
\]
and
\[
    \norm{T[X]}^2=\Theta(K),
    \qquad
    \gamma_{\mathrm{loc}}=O(K^{-1}).
\]
It uses
\[
    n=m+1=\frac{K^4}{8}+1=\Theta(\kappa^2).
\]
Thus it rules out estimates uniform over the dimension, but it is not a
fixed-dimensional counterexample. Whether an analogous obstruction exists
for fixed $n$, or with $n=o(\kappa^2)$, is not resolved by this construction.
\end{remark}

\begin{remark}[Where Hessian compatibility enters]
The direct computation in \Cref{rem:crosscheck} shows that the magnitude of
$T[X]$ is already encoded in the $t$-dependence of the $yy$ Hessian block.
Codazzi compatibility supplies an alternative derivation of the same
identity and, more importantly, yields the polarized second-Hermite identity
\eqref{eq:second-Hermite-id}. The latter makes $H$ and the packed local
generator self-adjoint. Changing only the $ty$ block while fixing the $yy$
block cannot change $T[X]$ for the test matrix used here.
\end{remark}

\begin{remark}[Minimum regularity]
The first- and second-Hermite identities use only commutation of
distributional derivatives. Thus $V\in C^2$ together with the bounded-Hessian
condition \eqref{eq:H1} is sufficient; no classical third derivative is
required. The tent-profile potential constructed above is $C^{2,1}$ and has
a globally Lipschitz Hessian, while \Cref{cor:smooth-ridge} gives a
$C^\infty$ realization.
\end{remark}

\begin{remark}[No circularity]
The local-gap upper bound uses only the explicit Rayleigh vector
\eqref{eq:Rayleigh-vector}, the independently established asymptotics of
$h_X$ and $\tau_X$, and Rayleigh--Ritz. It does not invoke a lower
spectral-gap estimate, a third-derivative moment bound, or a logarithmic
coupling inequality. A general lower bound may therefore be applied
afterward to identify order-sharpness without entering the proof above.
\end{remark}

% INTEGRATION-READY COROLLARY (intentionally inactive in this standalone note):
% After importing the newer local spectral theorem with
% Gamma=min{beta_* - 1, sqrt(n)(4sqrt(2t_0)+4t_0+4)}, the present family has
% Gamma=beta_* - 1=Theta(K). Its Schur bound gives tau_H^2<=1+Gamma, while
% Lemma~\ref{lem:TX-size} gives tau_H^2>=K/2-O(1); hence
% tau_H^2=Theta(K)=Theta(kappa_*^{1/2}). Combining the imported lower bound
% gamma_*>=max{alpha_*,1/(4+Gamma)} with Lemma~\ref{lem:Rayleigh} gives
% gamma_*=Theta(K^{-1})=Theta(kappa_*^{-1/2}). This saturates the
% beta_*-1 branch, not the sqrt(n) log(kappa_*) branch, and says nothing about
% fixed dimension.
\end{document}
